\documentclass[9pt,letterpaper]{article}
\usepackage[utf8]{inputenc}
\usepackage[T1]{fontenc}
\usepackage[spanish,es-nodecimaldot]{babel}
\usepackage[margin=0.55in]{geometry}
\usepackage{enumitem}
\usepackage[hidelinks]{hyperref}
\usepackage{titlesec}
\usepackage{microtype}
\usepackage{lmodern}
\pagestyle{empty}
\setlength{\parindent}{0pt}
\setlength{\parskip}{0pt}
\urlstyle{same}
\setlist[itemize]{leftmargin=1.2em,label=\textbullet,itemsep=0pt,topsep=1pt,parsep=0pt,partopsep=0pt}
\titleformat{\section}{\normalsize\bfseries}{}{0pt}{\MakeUppercase}[\vspace{-4pt}\rule{\textwidth}{0.45pt}]
\titlespacing*{\section}{0pt}{5pt}{2.5pt}
\newcommand{\resumeSubheading}[4]{\textbf{#1} \hfill #2\\[-1pt]\textit{#3} \hfill \textit{#4}\\[-1pt]}
\newcommand{\resumeProject}[2]{\textbf{#1} \hfill \textit{#2}\\[-1pt]}
\begin{document}
\begin{center}
{\Large\bfseries Diego Fernando Bustos Horta}\\[3pt]
Bogotá, Colombia \;|\; +57 310 610 5923 \;|\; \href{mailto:dbustoshu@unal.edu.co}{dbustoshu@unal.edu.co} \;|\; \href{https://www.linkedin.com/in/diegobstos/}{LinkedIn}
\end{center}
\section{Perfil}
Soy ingeniero mecánico en formación en la Universidad Nacional de Colombia. Me interesa el análisis de esfuerzos y la flexibilidad de tuberías, el comportamiento de los materiales y el diseño de sistemas mecánicos, y me gustaría orientar mi carrera hacia proyectos del sector oil \& gas. He trabajado con CAESAR II y la norma ASME B31.3 en casos prácticos, además de ensayos de laboratorio y modelamiento computacional.
\section{Educación}
\resumeSubheading{Universidad Nacional de Colombia -- Sede Bogotá}{Bogotá, Colombia}{Ingeniería Mecánica}{En curso}
\section{Formación complementaria}
\begin{itemize}
\item \textbf{Introducción al análisis de esfuerzos y deformaciones en tuberías -- casos de estudio prácticos (20 h)} --- IDOM / Universidad Nacional de Colombia, finalizado. Análisis de flexibilidad de sistemas de succión y descarga de bombas de crudo en CAESAR II 2025 bajo ASME B31.3: esfuerzos sostenidos, expansión térmica y prueba hidrostática; verificación de cargas en boquillas de tanques y bombas; cálculo de espesor por presión interna; diagnóstico de fallas y recomendaciones de rediseño de soportes.
\item \textbf{Machine Learning en Python y R} --- Udemy, en curso.
\item \textbf{SQL y Bases de Datos Relacionales} --- Udemy, en curso.
\item \textbf{Biomateriales y aplicaciones en salud} --- Universidad Nacional de Colombia.
\end{itemize}
\section{Proyectos académicos y de investigación}
\resumeProject{Modelamiento y control de una planta térmica}{Fundamentos de Control}
\begin{itemize}
\item Implementación y caracterización experimental de una planta térmica mediante Arduino.
\item Modelamiento del sistema y diseño de un controlador para regular su comportamiento.
\end{itemize}
\resumeProject{Participación en el Semillero GRISEC (energías renovables)}{Semillero GRISEC}
\begin{itemize}
\item Participación en trabajo experimental de caracterización de materiales: ensayos de tensión, flexión y dureza, así como tratamientos térmicos.
\item Análisis de resultados experimentales y estudio de propiedades de materiales y compuestos.
\end{itemize}
\resumeProject{Diseño de hidrolavadora de agua caliente}{Proyecto aplicado de ingeniería | En desarrollo}
\begin{itemize}
\item Diseño y dimensionamiento de sistemas térmicos, hidráulicos, estructurales y eléctricos.
\item Evaluación de componentes y alternativas de funcionamiento mediante cálculos y criterios de diseño.
\end{itemize}
\resumeProject{Diseño de máquina para extrusión y peletizado de plástico reciclado}{Universidad Nacional de Colombia | 2025}
\begin{itemize}
\item Diseño y dimensionamiento de un sistema de extrusión de polipropileno reciclado.
\item Selección y cálculo de transmisión mecánica, motor y componentes según potencia, torque y capacidad de producción.
\end{itemize}
\section{Experiencia}
\resumeSubheading{Universidad Nacional de Colombia -- Sede Bogotá}{Bogotá, Colombia}{Becario -- Ciencia e Ingeniería de los Materiales}{Sep. 2026 -- Dic. 2026}
\begin{itemize}
\item Ensayos de esfuerzos residuales por el método del agujero ciego en placas metálicas, bajo la dirección de la profesora Mónica Monsalve.
\item Apoyo a la gestión académica de la asignatura.
\end{itemize}
\resumeSubheading{Fisioestetic}{Bogotá, Colombia}{Gestión administrativa y operativa}{Ago. 2025 -- Jun. 2026}
\begin{itemize}
\item Desarrollo de herramientas en Excel para la gestión y análisis de inventarios, ventas y gastos operativos.
\item Elaboración de reportes e indicadores para el seguimiento de resultados y apoyo a la toma de decisiones.
\item Análisis de ventas, consumo y abastecimiento para identificar oportunidades de mejora en los procesos operativos.
\end{itemize}
\section{Habilidades técnicas}
\textbf{Análisis de tuberías y normatividad:} CAESAR II, ASME B31.3, análisis de flexibilidad y esfuerzos.
\textbf{Programación y análisis de datos:} Python, R, SQL, Power BI y Excel.
\textbf{Diseño y simulación:} Autodesk Inventor, AutoCAD, ANSYS y MATLAB.
\textbf{Control e instrumentación:} Arduino, control PI, PWM, modelamiento e identificación de sistemas.
\textbf{Idiomas:} Español (nativo); Inglés (intermedio, B1).
\textbf{Competencias:} Pensamiento analítico, resolución de problemas, investigación, aprendizaje autónomo y comunicación técnica.
\end{document}
